\section{Многочлен Тейлора и формула Тейлора с остаточным членом в форме Пеано. Единственность представления.}

\subsection{Многочлен Тейлора}

\begin{definition}
    Пусть функция \(f(x)\) имеет \(n\) производных в точке \(x_0\). 
    Многочленом Тейлора степени \(n\) называется
    \[
    T_n(x) = \sum_{k=0}^n \frac{f^{(k)}(x_0)}{k!} (x - x_0)^k.
    \]
    При \(x_0 = 0\) этот многочлен называют также многочленом Маклорена.
\end{definition}

\begin{theorem}[Свойство многочлена Тейлора]
    Для любого \(k = 0, 1, \dots, n\) выполнено
    \[
    T_n^{(k)}(x_0) = f^{(k)}(x_0).
    \]
\end{theorem}

\begin{Proof}
    При дифференцировании \(T_n(x)\) слагаемые со степенью меньше \(k\) обращаются в ноль (производная константы), 
    слагаемые со степенью больше \(k\) содержат множитель \((x - x_0)\) и при подстановке \(x = x_0\) дают ноль. 
    Остаётся только слагаемое с \(k\)-й степенью:
    \[
    \left( \frac{f^{(k)}(x_0)}{k!} (x - x_0)^k \right)^{(k)} = \frac{f^{(k)}(x_0)}{k!} \cdot k! = f^{(k)}(x_0).
    \]
\end{Proof}

\subsection{Формула Тейлора с остаточным членом в форме Пеано}

\begin{theorem}[Локальная формула Тейлора]
    Если существует \(f^{(n)}(x_0)\) (т.е. функция \(n\) раз дифференцируема в точке \(x_0\)), то
    \[
    f(x) = T_n(x) + o((x - x_0)^n) \quad \text{при } x \to x_0,
    \]
    где \(T_n(x)\) — многочлен Тейлора, а \(o((x - x_0)^n)\) обозначает функцию, для которой
    \[
    \lim_{x \to x_0} \frac{o((x - x_0)^n)}{(x - x_0)^n} = 0.
    \]
\end{theorem}

\begin{Proof}
    Рассмотрим \(R_n(x) = f(x) - T_n(x)\). Нам нужно доказать, что
    \[
    \lim_{x \to x_0} \frac{R_n(x)}{(x - x_0)^n} = 0.
    \]
    
    Применим правило Лопиталя \(n-1\) раз. Заметим, что
    \[
    R_n(x_0) = R_n'(x_0) = \dots = R_n^{(n-1)}(x_0) = 0,
    \]
    так как значения функции и её производных в точке \(x_0\) совпадают с многочленом Тейлора.
    
    Тогда:
    \[
    \lim_{x \to x_0} \frac{R_n(x)}{(x - x_0)^n} = 
    \lim_{x \to x_0} \frac{R_n'(x)}{n(x - x_0)^{n-1}} = 
    \dots = 
    \lim_{x \to x_0} \frac{R_n^{(n-1)}(x)}{n!(x - x_0)}.
    \]
    
    На этом шаге правило Лопиталя применить нельзя, так как \(f^{(n-1)}\) может не быть дифференцируемой в окрестности точки \(x_0\) (из условия следует только дифференцируемость в самой точке).
    
    Однако \(f^{(n-1)}\) дифференцируема в точке \(x_0\). Запишем для неё определение производной:
    \[
    f^{(n-1)}(x) = f^{(n-1)}(x_0) + f^{(n)}(x_0)(x - x_0) + o(x - x_0).
    \]
    Аналогично для многочлена Тейлора:
    \[
    T_n^{(n-1)}(x) = T_n^{(n-1)}(x_0) + T_n^{(n)}(x_0)(x - x_0) + o(x - x_0).
    \]
    
    Но \(f^{(n-1)}(x_0) = T_n^{(n-1)}(x_0)\) и \(f^{(n)}(x_0) = T_n^{(n)}(x_0)\) по свойству многочлена Тейлора. Вычитая, получаем:
    \[
    R_n^{(n-1)}(x) = f^{(n-1)}(x) - T_n^{(n-1)}(x) = o(x - x_0) - o(x - x_0) = o(x - x_0).
    \]
    
    Следовательно,
    \[
    \lim_{x \to x_0} \frac{R_n^{(n-1)}(x)}{n!(x - x_0)} = \lim_{x \to x_0} \frac{o(x - x_0)}{n!(x - x_0)} = 0.
    \]
    
    Таким образом, \(\displaystyle \lim_{x \to x_0} \frac{R_n(x)}{(x - x_0)^n} = 0\), что и означает \(R_n(x) = o((x - x_0)^n)\).
\end{Proof}

\subsection{Единственность представления}

\begin{theorem}[О единственности]
    Пусть \(\exists f^{(n)}(x_0)\) и 
    \[
    f(x) = P_n(x) + o((x - x_0)^n), \quad x \to x_0,
    \]
    где \(P_n(x)\) — многочлен степени не выше \(n\). Тогда \(P_n(x) = T_n(x)\).
\end{theorem}

\begin{Proof}
    По условию имеем два представления:
    \[
    f(x) = T_n(x) + o((x - x_0)^n), \quad
    f(x) = P_n(x) + o((x - x_0)^n).
    \]
    Вычитая, получаем:
    \[
    T_n(x) - P_n(x) = o((x - x_0)^n).
    \]
    
    Запишем \(P_n(x)\) в виде 
    \[
    P_n(x) = \sum_{k=0}^n a_k (x - x_0)^k.
    \]
    Тогда
    \[
    T_n(x) - P_n(x) = \sum_{k=0}^n \left( \frac{f^{(k)}(x_0)}{k!} - a_k \right) (x - x_0)^k = o((x - x_0)^n).
    \]
    
    Перейдем к пределу при \(x \to x_0\). Так как все члены суммы, кроме \(k = 0\), содержат множитель \((x - x_0)\) и стремятся к 0, то
    \[
    \lim_{x \to x_0} (T_n(x) - P_n(x)) = \frac{f(x_0)}{0!} - a_0 = 0 \implies a_0 = f(x_0).
    \]
    
    Разделим равенство на \((x - x_0)\):
    \[
    \frac{T_n(x) - P_n(x)}{x - x_0} = \left( \frac{f'(x_0)}{1!} - a_1 \right) + \sum_{k=2}^n \left( \frac{f^{(k)}(x_0)}{k!} - a_k \right) (x - x_0)^{k-1} = o((x - x_0)^{n-1}).
    \]
    Переходя к пределу при \(x \to x_0\), получаем \(a_1 = f'(x_0)\).
    
    Продолжая этот процесс, последовательно получим
    \[
    a_k = \frac{f^{(k)}(x_0)}{k!} \quad \text{для всех } k = 0, 1, \dots, n.
    \]
    Следовательно, \(P_n(x) = T_n(x)\).
\end{Proof}
